\section{Managed services: databases, RPC, archive, and indexing}
\label{sec:managed-services}

A stateless application on \Flux tolerates a bad host: if the container dies, \FluxOS
reschedules it elsewhere and the workload resumes with no state to lose. A managed
database cannot make that trade. Its correctness is defined by what happens to data
that exists only on machines the operator does not control, that are geographically
dispersed, that may leave the network without notice, and that share none of a
datacenter's usual assumptions --- no shared storage fabric, no operator-trusted
network, no guaranteed clock, no assurance that the node reporting a given IP is the
node that was there a minute ago. Every other node-level daemon this paper describes
(telemetry shipping, coordinated shutdown, expiry notification) is stateless with
respect to the application: it observes or forwards, and if it fails the application's
own state is untouched. A managed database is the one place on \Flux where the
platform's usual escape hatch --- reschedule and move on --- is exactly the operation
that can destroy data. This section describes the two implementations that exist for
this problem today, states precisely what each one guarantees, and shows that the two
guarantees are not the same guarantee.

\subsection{\texttt{Flux-Shared-DB}: ad-hoc leader election over the \FluxOS location API}
\label{subsec:flux-shared-db}

\texttt{Flux-Shared-DB} (\shipped) is a Node.js ``Operator'' sidecar that clusters
independent MySQL/MariaDB engine instances running on separate \FluxNode s into one
logical database. It intercepts the MySQL wire protocol, sequences writes on an
elected master, and replays them on followers via its own replication log, called the
``backlog'' \srcdoc{Flux-Shared-DB/README.md:5}. It is deployed as an ordinary
tenant-deployable \FluxOS application component \srcdoc{Flux-Shared-DB/README.md:20-49}.

\paragraph{Leader election.} Membership and master selection are derived entirely from
the \FluxOS application-location API, not from a consensus protocol: there are no term
numbers, no quorum proof, and no Raft or Paxos implementation anywhere in the
repository. \texttt{findMaster()} fetches the app's instance IPs from
\texttt{https://api.runonflux.io/apps/location/\{appName\}}, polls every reachable
peer's status, and ranks candidates by (1) highest replication sequence number,
descending; (2) a static IP, true first; (3) OS uptime, descending
\srccode{Flux-Shared-DB/\allowbreak{}ClusterOperator/\allowbreak{}Operator.js}{1464-1489}. If the local node
ranks first, it does not declare itself master outright: it asks the second-ranked
candidate to confirm before doing so. If the local node does not rank first, it asks
the top-ranked candidate who the master is, and if that answer disagrees with the
candidate list, it asks the \emph{reported} master to confirm itself
\srccode{Flux-Shared-DB/\allowbreak{}ClusterOperator/\allowbreak{}Operator.js}{1494-1543}. Any ambiguous or null
response causes the routine to recurse with no fixed retry cap or backoff
\srccode{Flux-Shared-DB/\allowbreak{}ClusterOperator/\allowbreak{}Operator.js}{1508,1519,1538,1552}. This is a
self-built, ad-hoc leader-election protocol layered directly on an HTTP location
service, and it is the default, unconditional path --- there is no feature flag
selecting a different mechanism.

\paragraph{Failover detection.} A health-check timer runs every $120{,}000$~ms
(2~minutes) and re-derives the master if any of: more than one node reports a
conflicting master IP; the master fails three consecutive \texttt{getMaster} probes;
the master's IP disappears from the app's location list; or, on the master itself, the
Socket.IO server has zero connected clients
\srccode{Flux-Shared-DB/\allowbreak{}ClusterOperator/\allowbreak{}Operator.js}{1265-1338}
\srccode{Flux-Shared-DB/\allowbreak{}ClusterOperator/\allowbreak{}server.js}{1543-1544}. Independently of the
poll, a follower's WebSocket \texttt{disconnect}\slash\texttt{connect\_error} handler
triggers immediate re-election on that follower, so the effective failover-detection
latency is bounded above by the 2-minute poll but can be much shorter when the
transport itself notices the drop
\srccode{Flux-Shared-DB/\allowbreak{}ClusterOperator/\allowbreak{}Operator.js}{178-196}.

\paragraph{Write durability.} This is the load-bearing finding for the section.

\begin{property}[\texttt{Flux-Shared-DB} acknowledges before it replicates]
\label{prop:sharedb-durability}
A write reaches the connecting MySQL client as soon as the elected master finishes
executing that query against the tenant's own database on the master's local engine
\srccode{Flux-Shared-DB/\allowbreak{}ClusterOperator/\allowbreak{}Backlog.js}{145-149}. Only \emph{after} the
client has been told the write succeeded does the master broadcast it to connected
follower sockets, via \texttt{serverSocket.emit('query', \ldots)}; no follower
acknowledgment is awaited or required before that call returns
\srccode{Flux-Shared-DB/\allowbreak{}ClusterOperator/\allowbreak{}Operator.js}{379-389}. No code path awaiting
a follower ack was found anywhere in \texttt{Operator.js}, \texttt{Backlog.js}, or
\texttt{server.js}: this is fire-and-forget asynchronous replication, and there is no
synchronous mode in this codebase. If the master fails between acknowledging the
client and completing that broadcast, the write is gone: the next master is whichever
surviving node holds the highest \emph{already-replicated} sequence number
\srccode{Flux-Shared-DB/\allowbreak{}ClusterOperator/\allowbreak{}Operator.js}{1483-1489}, and a write that never
left the old master leaves no trace anywhere else in the code reviewed. The size of
this data-loss window is not bounded by anything in the codebase examined: in the best
case it is the single most recent write; under sustained load or follower lag it can be
an unbounded number of recent writes, limited only by how far behind the surviving
followers happened to be.
\end{property}

A second, compounding detail: even the master's own record of the write in its local
backlog table is not confirmed before the client is told the write succeeded. The
insert into \texttt{flux\_backlog.\allowbreak{}backlog} is issued without an \texttt{await}, and the
subsequent error check on its result is short-circuited by a literal
\texttt{if (false \&\& \ldots)} guard that can never fire
\srccode{Flux-Shared-DB/\allowbreak{}ClusterOperator/\allowbreak{}Backlog.js}{126-144}. Reads compound the
problem in the other direction: ordinary (non-write) queries execute directly against
whichever node the client happens to be connected to, regardless of replication lag,
so a client on a lagging follower can read stale data indefinitely and two clients on
two different followers can observe different states at the same instant
\srccode{Flux-Shared-DB/\allowbreak{}ClusterOperator/\allowbreak{}Operator.js}{720-733}. An
\texttt{AUTH\_MASTER\_ONLY} setting restricts DB connections to the master only, but it
is opt-in and defaults to \texttt{false} \srccode{Flux-Shared-DB/\allowbreak{}ClusterOperator/\allowbreak{}config.js}{24}.

\paragraph{Partition behavior.} There is no quorum requirement to elect a master. A
minority partition containing the node with the highest sequence number will elect
that node as master regardless of how much of the mesh it can actually reach, so a
classic split-brain --- two nodes simultaneously believing themselves master and
accepting writes --- is possible whenever the mesh disagrees about who the master is
for up to the 2-minute health-check interval; the code's only defense is detecting
that disagreement on the next cycle, not preventing it in the first place
\srccode{Flux-Shared-DB/\allowbreak{}ClusterOperator/\allowbreak{}Operator.js}{1265-1271}. No explicit
split-brain test or fencing mechanism was found in the repository, so this specific
scenario is stated here as a structural consequence of the mechanisms above, not as an
observed incident \srcinferred.

\paragraph{Maturity, honestly assessed.} \texttt{Flux-Shared-DB} is shipped: it is the
default, unconditional replication path, and a CI pipeline builds and pushes three
Docker Hub image variants (\texttt{latest}, \texttt{latest-mysql},
\texttt{latest-mariadb}) on every merge to \texttt{master}
\srccode{Flux-Shared-DB/\allowbreak{}.github/\allowbreak{}workflows/\allowbreak{}docker-image.yml}{1-33}, which is evidence of
active use, not merely of a working build. But none of the properties above --- leader
election correctness, split-brain avoidance, or the failover data-loss window in
Property~\ref{prop:sharedb-durability} --- has an automated test. The only test file in
the repository exercises Bitcoin-message-signature verification, unrelated to
replication, failover, or backup/restore
\srccode{Flux-Shared-DB/\allowbreak{}test/\allowbreak{}bitcoinMessage.test.js}{1-18}. The hardest properties of
this system are, in this precise sense, unverified.

Two further findings, stated as facts rather than as speculation about intent: (1) an
AES-256-CBC ``comm'' encryption layer for the internal cluster protocol exists in code
(\texttt{Security.\allowbreak{}encryptComm}, a \texttt{shareKeys} socket event), but the
key-exchange handshake that would populate the AES key and IV on the connecting side is
commented out, so this layer is not exercised on the path reviewed and internal cluster
traffic --- queries, keys, session tokens --- runs as plaintext Socket.IO by default
\srccode{Flux-Shared-DB/\allowbreak{}ClusterOperator/\allowbreak{}Operator.js}{151-168}. (2) The MySQL wire-protocol
proxy parses the client's username and password scramble at connection time but never
verifies the scramble against a password
\srccode{Flux-Shared-DB/\allowbreak{}lib/\allowbreak{}mysqlServer.js}{291-320}; the operator's actual
authorization gate is an IP allowlist (\texttt{WHITELIST}, the first \texttt{172.x}
address to connect, or membership in the app's own instance-IP list)
\srccode{Flux-Shared-DB/\allowbreak{}ClusterOperator/\allowbreak{}Operator.js}{324-359}. Authentication to this
managed database is, in the code reviewed, IP-allowlist-only.

\subsection{\texttt{flux-\allowbreak{}pg-\allowbreak{}cluster}: Patroni/etcd-backed PostgreSQL HA}
\label{subsec:flux-pg-cluster}

\texttt{flux-\allowbreak{}pg-\allowbreak{}cluster} (\shipped) is a Go agent (\texttt{flux-agent}) that turns three
or more \FluxOS application instances into a standard Patroni-managed PostgreSQL
high-availability cluster, adding \Flux-specific glue: membership discovery via the
\FluxOS location API, a primary-routing TCP proxy, and conservative bootstrap/recovery
logic intended to prevent split-brain across an unreliable, geographically dispersed,
possibly NAT'd node set \srcdoc{flux-pg-cluster/README.md:7}
\srccode{flux-pg-cluster/\allowbreak{}cmd/\allowbreak{}flux-agent/\allowbreak{}daemon.go}{25}.

\paragraph{Leader election: standard, not custom.} Unlike \texttt{Flux-Shared-DB}, this
system does not reimplement consensus. \texttt{etcd} holds the distributed
configuration store, and Patroni elects the primary via \texttt{etcd} leader keys with
\texttt{PATRONI\_TTL}$=30$\,s and \texttt{PATRONI\_LOOP\_WAIT}$=10$\,s, both defaults
\srccode{flux-pg-cluster/\allowbreak{}patroni.yml.tpl}{30-33}
\srccode{flux-pg-cluster/\allowbreak{}internal/\allowbreak{}config/\allowbreak{}config.go}{150-151}. The \Flux-specific code
renders Patroni's configuration template from discovered peer IPs and manages the
\texttt{etcd} process lifecycle and cluster-membership reconciliation; it does not
reimplement the election itself \srccode{flux-pg-cluster/\allowbreak{}cmd/\allowbreak{}flux-agent/\allowbreak{}init.go}{422}.

\paragraph{Primary-routing proxy.} A Go TCP proxy listening on port 5433 polls every
candidate node's Patroni \texttt{/primary} and \texttt{/replica} REST endpoints every
\texttt{PROXY\_HEALTH\_INTERVAL\_SECONDS}$=3$\,s (default) and forwards new TCP
connections to whichever node most recently answered 200 on \texttt{/primary}
\srccode{flux-pg-cluster/\allowbreak{}cmd/\allowbreak{}flux-agent/\allowbreak{}proxy.go}{64,114-154}
\srccode{flux-pg-cluster/\allowbreak{}internal/\allowbreak{}patroni/\allowbreak{}client.go}{49-57}. In-flight connections are
not killed on failover --- the documentation states this explicitly, and the code
performs no active disconnect \srccode{flux-pg-cluster/\allowbreak{}cmd/\allowbreak{}flux-agent/\allowbreak{}proxy.go}{26-27}
--- so an application holding an open connection across a failover must detect the stale
connection itself.

\paragraph{Bootstrap and recovery gating.} On a from-scratch or dead-DCS boot, the
agent probes every expected peer's HMAC-authenticated identity endpoint and requires
either (a) exactly one node reporting non-empty PostgreSQL data, or (b) unanimous
total-DCS-loss with matching PostgreSQL system IDs across every reachable peer, before
selecting a recovery authority; ties are broken by the most advanced replica's durable
WAL position, with deterministic node-name tie-breaking as a last resort
\srccode{flux-pg-cluster/\allowbreak{}cmd/\allowbreak{}flux-agent/\allowbreak{}bootstrap\_identity.go}{263-378}. Any
unreachable peer, membership mismatch, or multiple-primary condition blocks automatic
recovery outright \srcdoc{flux-pg-cluster/README.md:206}
\srccode{flux-pg-cluster/\allowbreak{}cmd/\allowbreak{}flux-agent/\allowbreak{}bootstrap\_identity.go}{271-368}. Separately,
a quorum-loss recovery path writes to a quorum-probe \texttt{etcd} key every reconcile
cycle; if that write fails for \texttt{Desired\allowbreak{}State\allowbreak{}Stability\allowbreak{}Cycles}$=3$ (default)
consecutive cycles while the local node still holds PostgreSQL data, the agent rebuilds
the \texttt{etcd} cluster with \texttt{-\allowbreak{}-\allowbreak{}force-\allowbreak{}new-\allowbreak{}cluster}, gated by a cooldown
\srccode{flux-pg-cluster/\allowbreak{}cmd/\allowbreak{}flux-agent/\allowbreak{}daemon.go}{121-153,388-427}. This path is
exercised end-to-end by \texttt{test\_ec3\_majority\_replacement.\allowbreak{}py}, which kills two of
three nodes and asserts recovery within a 300-second test budget
\srccode{flux-pg-cluster/\allowbreak{}tests/\allowbreak{}scenarios/\allowbreak{}test\_ec3\_majority\_replacement.py}{1-60}.
\texttt{ALLOW\_PG\_DATA\_WIPE} defaults to \texttt{false}; on a system-ID mismatch, data
is wiped only if this flag is explicitly \texttt{true} \emph{and} the live primary
confirms authority, otherwise the agent logs an error and leaves the data directory
intact \srccode{flux-pg-cluster/\allowbreak{}internal/\allowbreak{}config/\allowbreak{}config.go}{137-141}.

\begin{remark}[Default replication favors availability over guaranteed durability]
\label{rem:pgcluster-durability}
Default replication in \texttt{flux-\allowbreak{}pg-\allowbreak{}cluster} is asynchronous: ``the primary commits
writes without waiting for replicas'' \srcdoc{flux-pg-cluster/README.md:210},
confirmed by \texttt{Patroni\allowbreak{}Synchronous\allowbreak{}Mode} defaulting to \texttt{false}
\srccode{flux-pg-cluster/\allowbreak{}internal/\allowbreak{}config/\allowbreak{}config.go}{158} and templated directly into
Patroni's \texttt{bootstrap.dcs} configuration
\srccode{flux-pg-cluster/\allowbreak{}patroni.yml.tpl}{37-39}. This is a documented data-loss-on-failover
trade-off, not an oversight, and the durability semantics themselves are delegated
entirely to upstream Patroni/PostgreSQL streaming replication rather than reimplemented
in this codebase. Setting \texttt{PATRONI\_SYNCHRONOUS\_MODE=\allowbreak{}true} is an opt-in toggle
that makes every commit wait for \texttt{PATRONI\_SYNCHRONOUS\_NODE\_COUNT}$=1$
(default) replica acknowledgment before returning success, eliminating the
failover-data-loss window at the cost of write latency; with
\texttt{PATRONI\_SYNCHRONOUS\_MODE\_STRICT=\allowbreak{}true} it additionally blocks all writes when
no synchronous replica is reachable \srcdoc{flux-pg-cluster/README.md:164-232} (the
config plumbing for this is confirmed in code
\srccode{flux-pg-cluster/\allowbreak{}internal/\allowbreak{}config/\allowbreak{}config.go}{158-161}; the runtime
commit-blocking behavior itself lives inside Patroni/PostgreSQL, outside this
repository's Go code, and is asserted here only as configuration wiring, not as
directly observed runtime behavior). Even in the default asynchronous mode,
\texttt{PATRONI\_MAX\_LAG}$=33{,}554{,}432$ bytes (32~MiB, default) caps how far behind
a replica may fall and still be eligible for promotion during failover
\srccode{flux-pg-cluster/\allowbreak{}internal/\allowbreak{}config/\allowbreak{}config.go}{153}
\srccode{flux-pg-cluster/\allowbreak{}patroni.yml.tpl}{34}: the tolerance is generous, but it is
bounded.
\end{remark}

\paragraph{Test coverage.} In sharp contrast to \texttt{Flux-Shared-DB}, this repository
has a pytest integration-test suite (\texttt{tests/\allowbreak{}scenarios/\allowbreak{}}) that spins up
three-node Docker Compose clusters and a mock \FluxOS location API and asserts on real
failure injection: leader kill followed by election within 120~s
(\texttt{test\_normal\_failover.\allowbreak{}py}), unreachable-node and late-API-registration
handling, majority-node replacement recovering via \texttt{force-\allowbreak{}new-\allowbreak{}cluster} within a
300-second budget (\texttt{test\_ec3\_majority\_replacement.\allowbreak{}py}), and system-ID
mismatch handling
\srccode{flux-pg-cluster/\allowbreak{}tests/\allowbreak{}scenarios/\allowbreak{}}{test\_normal\_failover.py, test\_ec3\_majority\_replacement.py, et al.}. It also has Go unit tests colocated
with the agent code (\texttt{bootstrap\_identity\_test.\allowbreak{}go}, \texttt{recovery\_test.\allowbreak{}go},
\texttt{proxy\_test.go}, \texttt{init\_test.go}, \texttt{config\_test.\allowbreak{}go},
\texttt{etcdmgr\_test.\allowbreak{}go}, \texttt{fluxapi/\allowbreak{}client\_test.\allowbreak{}go}). The repository's only CI
workflow builds and pushes Docker images and does not invoke the pytest suite, so the
suite is not a merge gate
\srccode{flux-pg-cluster/\allowbreak{}.github/\allowbreak{}workflows/\allowbreak{}docker-image.yml}{1-55}.
\texttt{DEAD\_CLUSTER\_RECOVERY} and \texttt{ALLOW\_PG\_DATA\_WIPE} are both live,
tested safety toggles gating destructive recovery paths (the former on by default, the
latter off by default)
\srccode{flux-pg-cluster/\allowbreak{}internal/\allowbreak{}config/\allowbreak{}config.go}{136-141} --- concrete evidence that
split-brain and data-loss edge cases were found and hardened against, rather than left
theoretical.

\subsection{The comparison}
\label{subsec:managed-db-comparison}

Both systems are \shipped, as operators a tenant deploys with its own application
(\S\ref{subsec:managed-db-tenant}); neither is the managed-database \emph{product} the public
roadmap lists as planned --- ``Postgres and MongoDB as one step in the deployment flow, with
backup and restore attached'' \srcdoc{flux-roadmap/public/flux-roadmap-public.md:71-72}. Both
solve the same problem: keeping a database consistent
across nodes that are untrusted, dispersed, and may vanish. The engineering postures
are not comparable.

\begin{center}
\small
\begin{tabular}{|p{3.6cm}|p{5.6cm}|p{5.6cm}|}
\hline
& \textbf{\texttt{Flux-Shared-DB}} & \textbf{\texttt{flux-\allowbreak{}pg-\allowbreak{}cluster}} \\
\hline
Election protocol & Ad-hoc, hand-rolled over the \FluxOS location API; no term numbers, no quorum proof \srccode{Flux-Shared-DB/\allowbreak{}ClusterOperator/\allowbreak{}Operator.js}{1464-1489} & Patroni/\texttt{etcd}, standard leader-lease election, TTL 30~s, \texttt{loop\_wait} 10~s \srccode{flux-pg-cluster/\allowbreak{}internal/\allowbreak{}config/\allowbreak{}config.go}{150-151} \\
\hline
Failover detection & 2-min health-check poll + immediate reconnect-triggered re-election \srccode{Flux-Shared-DB/\allowbreak{}ClusterOperator/\allowbreak{}server.js}{1543-1544} & Patroni lease expiry (TTL 30~s) + proxy health poll every 3~s \srccode{flux-pg-cluster/\allowbreak{}internal/\allowbreak{}config/\allowbreak{}config.go}{165} \\
\hline
Default write durability & Ack after master's local write only; broadcast to followers is unawaited (Property~\ref{prop:sharedb-durability}) & Ack after primary's local commit; async streaming replication (Remark~\ref{rem:pgcluster-durability}) \\
\hline
Synchronous option & None found in the codebase & \texttt{PATRONI\_SYNCHRONOUS\_MODE=\allowbreak{}true}, opt-in \srccode{flux-pg-cluster/\allowbreak{}internal/\allowbreak{}config/\allowbreak{}config.go}{158} \\
\hline
Split-brain defenses & Detection only, on next 2-min cycle; no fencing found \srccode{Flux-Shared-DB/\allowbreak{}ClusterOperator/\allowbreak{}Operator.js}{1265-1271} & Quorum-loss gating, unanimous-identity bootstrap check, \texttt{force-\allowbreak{}new-\allowbreak{}cluster} cooldown \srccode{flux-pg-cluster/\allowbreak{}cmd/\allowbreak{}flux-agent/\allowbreak{}daemon.go}{121-153,388-427} \\
\hline
Automated failure-injection tests & None (only a signature-verification unit test exists) \srccode{Flux-Shared-DB/\allowbreak{}test/\allowbreak{}bitcoinMessage.test.js}{1-18} & pytest suite killing nodes and asserting recovery within stated time budgets \srccode{flux-pg-cluster/\allowbreak{}tests/\allowbreak{}scenarios/\allowbreak{}test\_ec3\_majority\_replacement.py}{1-60} \\
\hline
\end{tabular}
\end{center}

The lesson is not that one team is better than the other; it is that on a decentralized
substrate, the durability guarantee a managed service can offer is bounded by the
replication protocol beneath it, and the two protocols here are not equivalent. A
tenant choosing between \texttt{Flux-Shared-DB} and \texttt{flux-\allowbreak{}pg-\allowbreak{}cluster} is choosing
between different, precisely different, durability guarantees --- not between two
skins on the same underlying service --- and the platform should say which is which
rather than presenting both under a single ``managed database'' label.

\begin{figure}[H]
\centering

\begin{tikzpicture}[
  font=\scriptsize, x=1mm, y=1mm,
  ax/.style={-{Latex[length=1.6mm]}, thick},
  ev/.style={circle, fill=black, inner sep=0.8pt},
  lbl/.style={font=\tiny, align=center, inner sep=1pt},
  nm/.style={font=\scriptsize, anchor=east},
  loss/.style={fill=black!14},
]
  % ---- timeline 1: Flux-Shared-DB
  \fill[loss] (18,34) rectangle (108,40);
  \draw[ax] (0,37) -- (112,37);
  \node[nm] at (-2,37) {\texttt{Flux-Shared-DB}};
  \node[ev] (w1)  at (6,37)  {}; \node[lbl, above=0.5mm of w1] {write};
  \node[ev] (l1)  at (18,37) {}; \node[lbl, below=0.5mm of l1] {local\\write};
  \node[ev] (a1)  at (18,37) {}; \node[lbl, above=0.5mm of a1] {\textbf{ack}};
  \node[ev] (b1)  at (46,37) {}; \node[lbl, below=0.5mm of b1] {broadcast\\emitted\\(unawaited)};
  %% y=37.2 put this on the axis itself; the rule struck through the text.
  \node[lbl, anchor=west] at (74,39.6) {follower receipt: not awaited, not bounded};

  % ---- timeline 2: flux-pg-cluster, async
  \fill[loss] (18,18) rectangle (58,24);
  \draw[ax] (0,21) -- (112,21);
  \node[nm] at (-2,21) {\texttt{flux-\allowbreak{}pg-\allowbreak{}cluster} (async)};
  \node[ev] (w2)  at (6,21)  {}; \node[lbl, above=0.5mm of w2] {write};
  \node[ev] (c2)  at (18,21) {}; \node[lbl, below=0.5mm of c2] {local\\commit};
  \node[lbl, above=0.5mm] at (18,21) {\textbf{ack}};
  \node[ev] (r2)  at (58,21) {}; \node[lbl, below=0.5mm of r2] {replica\\receives};
  \node[ev] (p2)  at (86,21) {}; \node[lbl, below=0.5mm of p2] {replica\\applies};

  % ---- timeline 3: flux-pg-cluster, synchronous
  \draw[ax] (0,5) -- (112,5);
  \node[nm] at (-2,5) {\texttt{flux-\allowbreak{}pg-\allowbreak{}cluster} (sync)};
  \node[ev] (w3)  at (6,5)  {}; \node[lbl, above=0.5mm of w3] {write};
  \node[ev] (c3)  at (18,5) {}; \node[lbl, below=0.5mm of c3] {local\\commit};
  \node[ev] (r3)  at (58,5) {}; \node[lbl, below=0.5mm of r3] {replica\\acknowledges};
  \node[lbl, above=0.5mm] at (58,5) {\textbf{ack}};
  \draw[very thick] (58,3.6) -- (58,6.4);
  %% Same as the first timeline: 0.6 above a rule at y=5 still crosses the glyphs.
  \node[lbl, anchor=west] at (64,7.6) {loss window collapses to zero};

  \node[lbl, anchor=west, fill=black!14, inner sep=1.2pt] at (0,46)
    {shaded: writes acknowledged to the client that a failover can still lose};
\end{tikzpicture}

\caption{Write-acknowledgment timing relative to replication, compared across
\texttt{Flux-Shared-DB} and the two \texttt{flux-\allowbreak{}pg-\allowbreak{}cluster} replication modes.}
\label{fig:managed-db-ack-timing}
\end{figure}

\subsection{How a tenant uses these}
\label{subsec:managed-db-tenant}

\paragraph{\texttt{Flux-Shared-DB}.} A tenant registers two \FluxOS application
components --- a stock \texttt{mysql:latest}\slash\texttt{mariadb} engine and the
\texttt{runonflux/\allowbreak{}shared-\allowbreak{}db} Operator image --- or a single bundled image with the
engine baked in \srcdoc{Flux-Shared-DB/README.md:20-49,73-101}. Credentials are set at
first boot via \texttt{DB\_INIT\_PASS}\slash\texttt{DB\_USER} environment variables and are
not re-keyable afterward: ``changing it later will not re-key an already-initialized
\texttt{/app/db} datadir'' \srcdoc{Flux-Shared-DB/README.md:88}. The tenant application
connects with an ordinary MySQL connection string to the Operator's proxy port (default
3307) \srcdoc{Flux-Shared-DB/README.md:46-49}. As established in
\S\ref{subsec:flux-shared-db}, authorization at that proxy is IP-allowlist-based, not
credential-based, and internal cluster traffic (queries, keys, session tokens) is
plaintext Socket.IO by default, since the AES ``comm'' layer's key exchange is
commented out. The debug UI can optionally run over TLS with a self-signed certificate
if \texttt{SSL=true}, default \texttt{false}
\srccode{Flux-Shared-DB/\allowbreak{}ClusterOperator/\allowbreak{}Security.js}{163-185}
\srccode{Flux-Shared-DB/\allowbreak{}ClusterOperator/\allowbreak{}config.js}{25}. No at-rest encryption of the
MySQL/MariaDB data directory was found in the repository.

\paragraph{\texttt{flux-\allowbreak{}pg-\allowbreak{}cluster}.} A tenant registers a single \FluxOS application
component with a fixed port set (5432, 5433, 8008, 2379, 2380)
\srcdoc{flux-pg-cluster/README.md:78}. Superuser and replication credentials are
supplied as \texttt{POSTGRES\_SUPERUSER\_PASSWORD}\slash\texttt{POSTGRES\_REPLICATION\_PASSWORD}
environment variables and persisted, shell-quoted, to \texttt{/\allowbreak{}etc/\allowbreak{}cluster\_env} for
reuse by later agent subcommands
\srccode{flux-pg-cluster/\allowbreak{}internal/\allowbreak{}config/\allowbreak{}config.go}{331-332}. Tenants connect through
the primary-routing proxy on port 5433, which always routes writes to the current
primary across failovers, or bypass it to a specific node on 5432 for read-only access
to a replica \srcdoc{flux-pg-cluster/README.md:284-338}. Encryption is opt-in:
\texttt{SSL\_ENABLED} defaults to \texttt{false}
\srccode{flux-pg-cluster/\allowbreak{}internal/\allowbreak{}config/\allowbreak{}config.go}{117}; when set, a deterministic
self-signed CA and per-service TLS certificates are generated for PostgreSQL,
\texttt{etcd} peer/client traffic, and the Patroni REST API, and \texttt{pg\_hba}
accepts \texttt{hostssl}\slash\texttt{cert clientcert=verify-\allowbreak{}full} for replication
connections alongside plain \texttt{host} md5 entries that the template and the
reconciler retain in either mode \srccode{flux-pg-cluster/\allowbreak{}generate-certs.sh}{45-306}
\srccode{flux-pg-cluster/\allowbreak{}patroni.yml.tpl}{69-73}
\srccode{flux-pg-cluster/\allowbreak{}cmd/\allowbreak{}flux-agent/\allowbreak{}daemon.go}{452-457}. With \texttt{SSL\_ENABLED=\allowbreak{}false}, the
default, inter-node PostgreSQL replication and \texttt{etcd} traffic run over the
public internet in plaintext --- the project's own architecture documentation labels
this path ``Replication via Public Internet'' \srcdoc{flux-pg-cluster/README.md:55}.
As with \texttt{Flux-Shared-DB}, no at-rest encryption of the PostgreSQL data directory
was found in the repository.

\subsection{Backup and restore}
\label{subsec:managed-db-backup}

\paragraph{\texttt{Flux-Shared-DB}.} Backups are periodic full \texttt{mysqldump}
snapshots, taken as part of the log-compaction cycle described in the evidence,
written to \texttt{./dumps/*.sql} inside the container and retaining only the two most
recent files \srccode{Flux-Shared-DB/\allowbreak{}ClusterOperator/\allowbreak{}Operator.js}{504-548}. If the
application's declared container data includes \texttt{s:/app/dumps}, \FluxOS's own
directory-sync mechanism replicates that path across the app's instances --- this
sync is a \FluxOS-level feature, not code in the operator repository itself
\srcdoc{Flux-Shared-DB/README.md:76,81-83}. Anyone who can reach the debug UI and
authenticate via a Bitcoin-message-signature login tied to the application owner's
ZelID can list, download, or delete backups, or execute an uploaded or
URL-fetched \texttt{.sql} file directly into the live cluster via
\texttt{/executebackup}, which first rolls the database back to sequence zero --- i.e.,
wipes it --- before importing \srccode{Flux-Shared-DB/\allowbreak{}ClusterOperator/\allowbreak{}server.js}{653-715}.
No evidence of restore being exercised by an automated test was found; the only test
file in the repository covers signature verification, unrelated to backup or restore
\srccode{Flux-Shared-DB/\allowbreak{}test/\allowbreak{}bitcoinMessage.test.js}{1-18}.

\paragraph{\texttt{flux-\allowbreak{}pg-\allowbreak{}cluster}.} An opt-in (\texttt{BACKUP\_ENABLED=\allowbreak{}false} by
default \srcdoc{flux-pg-cluster/README.md:170}) backup agent runs \texttt{pg\_dumpall},
gzip-compresses the result, and overwrites a previous backup only after an integrity
check (valid gzip, and presence of the ``PostgreSQL database dump complete'' marker)
via a temp-file-then-atomic-rename pattern
\srccode{flux-pg-cluster/\allowbreak{}cmd/\allowbreak{}flux-agent/\allowbreak{}backup.go}{90-143,204-223}. Backups run only on
whichever node Patroni currently reports as the healthy running primary, avoiding
duplicate copies
\srccode{flux-pg-cluster/\allowbreak{}cmd/\allowbreak{}flux-agent/\allowbreak{}backup.go}{93-99,145-164}. Retention keeps the
newest \texttt{BACKUP\_RETENTION\_COUNT}$=1$ (default) files, then prunes further by a
total-size cap that is unlimited by default
\srccode{flux-pg-cluster/\allowbreak{}cmd/\allowbreak{}flux-agent/\allowbreak{}backup.go}{275-337}. Backup storage is
\texttt{/\allowbreak{}var/\allowbreak{}lib/\allowbreak{}postgresql/\allowbreak{}backups}, explicitly documented as \emph{not} covered by
persisting the main data-directory volume, and lost on reschedule unless separately
persisted \srcdoc{flux-pg-cluster/README.md:77,276}. Restore is a manual, documented
\texttt{gunzip~|~psql} command \srcdoc{flux-pg-cluster/README.md:280-282}; the README
itself instructs the operator to ``verify restores regularly'' as their own
responsibility \srcdoc{flux-pg-cluster/README.md:186}, which is a stronger honesty
signal than a silent absence of guidance but does not amount to an automated restore
guarantee. A file named \texttt{test\_restore\_reconciliation.\allowbreak{}py} exists in the same
test suite \srccode{flux-pg-cluster/\allowbreak{}tests/\allowbreak{}scenarios/\allowbreak{}}{test\_restore\_reconciliation.py},
but its contents were not reviewed for this section, so it is not counted here as
verified automated restore coverage.

For neither system, then, has restore been exercised in a way this section can verify
end to end: \texttt{Flux-Shared-DB} has no restore test at all, and
\texttt{flux-\allowbreak{}pg-\allowbreak{}cluster}'s restore path is documented as a manual operation with an
unreviewed test file of unknown scope.

\subsection{RPC, archive, and indexing (\planned)}
\label{subsec:rpc-archive-indexing}

The roadmap names RPC, archive, and indexing services the network's ``first
capacity-selling product,'' scheduled for Q3 2026
\srcroadmap{flux-roadmap-public.md:33-34}, and states that this product ``pairs
directly with chain pruning: the archive nodes pruning requires become the product''
\srcroadmap{flux-roadmap-public.md:34}. Main-chain pruning itself is treated at
reimplementation depth in \S8; the roadmap's own framing of pruning is that ``pruned
nodes stay fully able to produce blocks, with designated archive nodes preserving full
history for explorers and the data services above''
\srcroadmap{flux-roadmap-public.md:121-122}. The direction of dependency between the two
items is stated loosely (``pairs directly with''), not as a strict blocking gate
\srcinferred.

No RPC, archive, or indexing implementation was surveyed for this section: the evidence
available covers only the two managed-database operators above, and no repository
matching an RPC gateway, an archive-node service, or an indexing service was read. This
item is therefore reported here purely as a roadmap commitment, and its construction is
left explicitly open rather than described as a system that does not yet exist.
\begin{remark}[The shape of the RPC product --- settled at the architectural level]
\label{rem:rpc-shape}
The request-routing model was carried as undecided. Its architecture is now settled, and in a way
that answers the structural half of the question while leaving the commercial half open:
\textbf{archive nodes run as \FluxCloud{} applications}, deployed by the team, with access to RPC and
other node APIs sold as a product \srcintent{08-answers-round-2.md, round 16} (\planned).

That decides three things. Routing is to \emph{operator-deployed shared endpoints} rather than to
per-tenant archive deployments --- a tenant buys API access, not a node. The pruning-exempt storage
requirement does not interact with the tiering and collateral model of \S\ref{sec:consensus} at all,
because an archive node is a tenant workload subject to the application resource model, not a
\FluxNode{} tier: it needs no collateral and confers no eligibility. And history availability is
consequently a property of this product rather than of the node set
(\Cref{rem:archive-product}).

What remains commercial rather than architectural: the indexing schema and the pricing unit --- per
request, per method, per byte returned, or a flat subscription. Those do not constrain any mechanism
in this paper and are not stated here.
\end{remark}

Whatever this product becomes, it will inherit the same question this section has just
answered for the two managed-database implementations: what does the platform promise
about durability and consistency, and under what replication or storage protocol.
Archived chain history is, if anything, a harder case than a tenant's own database,
because it is by definition unrecoverable once lost from every archive node
simultaneously; the paper takes no position here on which of the two postures compared
above --- hand-rolled and untested, or consensus-backed and failure-tested --- an
archive service would follow, because no such service has been surveyed.

\subsection{Implication for lifecycle: eviction is where the guarantee is tested}
\label{subsec:managed-db-lifecycle}

The containers that implement both systems above are, from the node's point of view,
ordinary Docker containers, and they are stopped through the same generic pipeline as
any other application when a node evicts, redeploys, or is asked to stop an
application. In production today, that pipeline executes only its Signal and Cleanup
stages; the drain and pre-stop hook stages exist in the state machine and are
unit-tested, but are ``not yet reachable from real events''
\srccode{flux-shutdownd/\allowbreak{}crates/\allowbreak{}shutdownd/\allowbreak{}src/\allowbreak{}pipeline/\allowbreak{}state\_machine.rs}{3-7}, even
though the underlying mechanism already reads per-container
\texttt{shutdown.\{drain,prestop,graceful\}-s} labels intended to drive exactly this
behavior \srccode{flux-shutdownd/\allowbreak{}crates/\allowbreak{}shutdownd/\allowbreak{}src/\allowbreak{}config.rs}{41-65}. For a
stateless workload, skipping a drain step is a latency blip. For a managed database, it
is the difference between an orderly checkpoint and whatever durability guarantee ---
or absence of one --- \S\ref{subsec:managed-db-comparison} describes: an application
that is killed rather than quiesced is, from the perspective of
Property~\ref{prop:sharedb-durability}, indistinguishable from a crash.

This is the requirement §16 must carry forward: a managed database is a stateful
service whose eviction destroys data, so the grace, suspension, and eviction state
machine that governs how and when a node stops an application matters more here than
for any stateless workload this paper describes elsewhere.

%% Drain this section's float queue before the next begins.
\FloatBarrier
